\documentclass[10pt,conference]{IEEEtran}
\usepackage{booktabs}
\usepackage{hyperref}
\usepackage{xurl}

\title{Evaluating Repository-Changing Agents Through State-Diff Contracts}
\author{\IEEEauthorblockN{Anonymous Author(s)}}

\begin{document}
\maketitle

\begin{abstract}
Repository-changing agents introduce a software maintenance problem: their
outputs are not only answers, but state transitions. A useful evaluation must
therefore ask which files, memory entries, review states, and tool-observable
facts changed, and whether those changes obeyed policy. This short paper
proposes state-diff contracts for evaluating such agents. A pilot artifact with
24 cases models tool calls, observations, permission gates, prompt-injection
conditions, memory scope, replay, and final-state assertions. The results are
not a general benchmark, but they show how state-diff and permission failures
can be surfaced as analyzable software artifacts. The contribution is a
software-analysis framing for maintaining agent workflows as prompts, models,
and tool schemas evolve.
\end{abstract}

\section{Introduction}

Software engineering agents increasingly operate on repositories and workflow
state. They edit files, inspect project context, run tests, prepare reports, and
request or bypass review. A final textual answer is insufficient evidence of
correct behavior. The maintainable unit is the trajectory: tool calls,
observations, policy checks, state diffs, and replay information.

This paper argues that repository-changing agents should be evaluated through
state-diff contracts. A state-diff contract specifies the initial state, allowed
actions, forbidden actions, and expected final state. It turns an agent run into
an analyzable artifact for software maintenance.

\section{Contract Model}

A contract has five parts:

\begin{enumerate}
\item task intent and trigger condition;
\item tool interface and allowed arguments;
\item permission policy for side effects;
\item state assertions over files, repository status, memory, and review state;
\item replay record sufficient to diagnose failure.
\end{enumerate}

The model is deliberately thinner than a full agent runtime. It does not impose
a planner. It inspects whether the external effects obey the contract.

\section{Pilot}

The anonymous artifact implements 24 cases across routing, stateful tool use,
permission and injection behavior, memory scoping, and replay/recovery. Three
prompt variants are compared: no harness, thick checklist, and thin contract.
The harness records invalid tool calls, permission violations, unsafe secret
access, state-diff correctness, parse failures, and replayability.

The pilot's purpose is to test the evaluation shape. The case set is too small
for general claims about providers. However, the traces demonstrate how failures
that are invisible in final answers become inspectable maintenance facts.

\section{Implications for Software Evolution}

Agent workflows evolve like software. Prompts change, tool schemas change,
model behavior changes, permission policies change, and downstream workflows
gain new side effects. Without replay and state-diff contracts, teams cannot
distinguish a model regression from a tool-schema bug, stale memory, overbroad
permission, or parser failure.

State-diff contracts provide a practical regression surface. They can be run
when a model is updated, a tool schema changes, a repository policy changes, or
a workflow gains a new side effect.

\section{Limitations}

The pilot is synthetic and small. The artifact uses mock state rather than
large-scale production repositories. The results are descriptive, not
inferential. The contribution is a tractable analysis framing and artifact
shape, not a comprehensive benchmark.

\section{Conclusion}

Repository-changing agents should be evaluated by what they do to software
state. State-diff contracts make those effects visible, testable, and replayable
as agent workflows evolve.

\begin{thebibliography}{9}
\bibitem{toolsandbox} ToolSandbox, stateful tool-use evaluation benchmark.
\bibitem{agentdojo} AgentDojo, dynamic tool-calling and prompt-injection evaluation.
\bibitem{sweagent} SWE-agent and agent-computer interfaces for software engineering agents.
\bibitem{deli} Deli AutoResearch framework for long-horizon autonomous tasks.
\end{thebibliography}

\end{document}

